\documentclass[10pt,a4paper]{amsart}
\usepackage[margin=19mm]{geometry}
\usepackage{amsmath,amssymb,mathtools}
\usepackage[scr=rsfs,cal=euler]{mathalfa}
\usepackage{xfrac}
\usepackage[unicode,hidelinks,bookmarks=false]{hyperref}
\newcommand{\ie}{i.e.\ }
\newcommand{\cf}{cf.\ }
\newcommand{\resp}{respectively\ }
\newcommand{\Subsection}{\subsection}
\providecommand{\nobreakdash}{\nobreak\hbox{-}}
\setlength{\emergencystretch}{2em}
\multlinegap0pt
\usepackage{mathtools}
\usepackage{dsfont}
\usepackage{tikz}
\usepackage{amssymb}
\newtheorem{lemma}{Lemma}
\newtheorem{proposition}{Proposition}
\newtheorem{theorem}{Theorem}
\usepackage{cleveref}
\crefname{lemma}{Lemma}{Lemmas}
\crefname{proposition}{Proposition}{Propositions}
\crefname{theorem}{Theorem}{Theorems}
\newcommand{\R}{\mathbb{R}}
\DeclareMathOperator{\clarke}{\partial^{*} \! \!}
\DeclareMathOperator{\dacy}{\delta_{\mathrm{Acy}}}
\DeclareMathOperator{\dhom}{d_{\mathrm{Hom}}}
\DeclareMathOperator{\diam}{\mathrm{diam}}
\DeclareMathOperator{\diff}{\mathrm{d} \!}
\newcommand{\eps}{\varepsilon}
\DeclareMathOperator{\lip}{\mathrm{Lip}}
\newcommand{\mres}{\mathbin{\vrule height 1.6ex depth 0pt width 0.13ex\vrule height 0.13ex depth 0pt width 1.3ex}}
\newcommand{\norme}[1]{\left\lvert \left \lvert #1 \right \rvert \right \rvert}
\title{Continuity of the normal cycle with respect to $C^0$-topologies}
\author{David Cohen-Steiner, Antoine Commaret}
\date{Source: arXiv:2609.05732v1 (CC BY 4.0)}
\begin{document}
\maketitle

\noindent\emph{Source and licence.} Continuity of the normal cycle with respect to $C^0$-topologies. Version arXiv:2609.05732v1 (CC BY 4.0), distributed by arXiv under the Creative Commons Attribution 4.0 International licence, http://creativecommons.org/licenses/by/4.0/, as recorded in the arXiv licence metadata for this exact version and shown on the abstract page https://arxiv.org/abs/2609.05732v1. Full licence notice: \texttt{licenses/arxiv-2609.05732-CC-BY-4.0.txt}.\par\smallskip

\noindent\textit{Editorial scope.} Counted unit: the article's theorem theorem (source0.tex lines 1097--1107). The article's complete original proof of it is delivered with it (source0.tex lines 1109--1127, one \texttt{proof} environment of 19 lines). No mathematics is written, repaired or re-derived: the statement and the proof are the article's own lines, in the article's own notation, and no retained line is substituted or re-flowed.  Retained uncounted as the source-local context the proof needs: lines 920--928; lines 972--979; lines 1013--1021; lines 1037--1045.  Nothing was omitted from any retained span, so the package carries no omission record.  No retained line contains a citation, so the wrapper carries no bibliography and no bibliography entry.  No retained line contains a figure environment, an image-inclusion command or drawing code, so no image is delivered and the package declares no asset dependency.  Editorial formatting only, in addition: an AMS \texttt{amsart} preamble that restates the article's own declarations and macros used by the retained lines, namely the environments \texttt{lemma}, \texttt{proposition}, \texttt{theorem} on separate counters, together with the standard packages those lines need.  The retained environments print in this document as Lemma 1, Proposition 1, Theorem 1 in order of appearance, whereas the article prints the same items with the numbering of its own full document.  The complete native source of the article as distributed in the arXiv e-print for this exact version is bundled unchanged as \texttt{sources/209461/source0.tex}.
    \begin{lemma}[Mass bound for Monge-Ampère currents, d.c case]
    \label{lem:ma_bound}
    For any function $h : \R^d \to \R$ of the form $h=f-g$ where $f,g$ are convex, we have for every compact subset $K$ of $\R^d$:
    \begin{equation*}
    \label{eq:mineurs}
        M([\diff h] \mres K \times \R^d)) \leq C_d (\diam(K) + 2d\max(\lip(f|_K),\lip(g|_K)))^d,
    \end{equation*}
    with $C_d = \frac{\omega_d}{d!}2^d \binom{2d}{d}$.
    \end{lemma}

    \begin{proposition}
    \label{prop:normal_bound}
    Let $h : \R^d \to \R_{\geq 0}$ be a d.c nondegenerate aura with $X \coloneqq h^{-1}(  0 )$ and let $\mu \coloneqq \lim_{t \to 0^+} \inf \{ \Delta (\clarke h(x) ) ,  0 < h(x) < t \}$. Denote $N_X$ the normal cycle of $X$. We have 
    \begin{equation*}
     M(N_X) \leq K_{d, \mu} M( [\diff h] \mres X \times \R^d ).
    \end{equation*}
    where $K_{d, \mu}$ is a constant depending on $d, \mu$.
    \end{proposition}

\begin{proposition}[Mass bound for normal cycles of sublevel sets of d.c functions]
\label{prop:sub_dc_aura}
Let $h = f-g: \R^d \to \R_{\geq 0}$ be a d.c function and let $t \in \R$ be such that $\mu = \lim_{\delta \to 0^+} \inf \{\Delta(\clarke \, h(x)), t < h(x) \leq t + \delta \}$ is positive. Then $X_t = h^{-1}(-\infty, t]$ is a WDC set and we have
\begin{equation*}
    M(N_{X_t}) \leq C_d K_{d, \mu}( \diam(X)+ 2d\max(\lip(f|_{X_t}), \lip(g|_{X_t})))^d,
\end{equation*}
where $C_d$ (resp. $K_{d, \mu}$) is the constant in \Cref{lem:ma_bound} (resp. \Cref{prop:normal_bound}). 
    
\end{proposition}

\begin{lemma}[Deformation retraction and smoothing]
\label{lem:flot_lisse}
   Let $h: \R^d \to \R_{\geq 0}$ be a $2L$-Lipschitz, nondegenerate aura and let $\rho, c$ be given as in the paragraph above. Then for every $\eps > 0$ and $0 < t \leq c$ there exists a positive $S(t, \eps)$, nondecreasing in $t$, such that for every $0 < s \leq \min(S(t, \eps), (c - t)/4L)$, 
   \begin{itemize}
       \item On $h_s^{-1}(t + 2sL)$, $\norme{\nabla h_s(x)} \geq \rho(1 - \eps)$.
       \item There exists a deformation retraction of $h_s^{-1}[0,t + 2sL]$ onto $X = h^{-1}(0)$ whose trajectories have length at most $(t + 4sL)/(\rho - \eps)$, so that  
       \[ \max( \dhom(X, h_s^{-1}[0,t + 2sL]), \dacy(X, h_s^{-1}[0,t + 2sL])) \leq (t + 4sL)/(\rho - \eps).\]
   \end{itemize}
\end{lemma}

\begin{theorem}
Let $h = f - g$ be a nondegenerate d.c aura such that $f$ and $g$ are $L$-Lipschitz on a neighborhood of $X = h^{-1}(0)$. Then for any positive decreasing sequence $t_n$ going to zero, there exists a positive sequence $s_n$ such that the sequence of compact domains $X_n \coloneqq h_{s_n}^{-1}[0, t_n + 2s_nL]$ satisfies the following for $n$ sufficiently large:
\begin{itemize}
\item $X_n$ is a smooth compact domain of $\R^d$,
\item $\limsup_n M(N_{X_n}) < + \infty$,
\item $X_{n+1} \subset X_n$,
\item $X = \bigcap_{n} X_n$,
\item $X_n$ converges to $X$ both in the acyclic convergence and with respect to the homotopy distance.
\end{itemize}

\end{theorem}

\begin{proof}

 Since $\norme{h_s - h_{s'}} \leq 2(s - s')L$, observe that when $t - t' \geq 2(s -s')L$, we have 
\begin{equation}
\label{eq:inclusion}
    h_{s'}^{-1}[0, t'] \subset h_s^{-1}[0,t]. 
\end{equation} 
Let $b_n = \min(S(t_1, \eps), \dots, S(t_n, \eps), 1/n)$ for some $0 < \eps < 1$ and $s_{k,n} = b_k + (t_n - t_k)/2L$. Note that $b_n$ is nonincreasing.
Put
\[ s_n = \inf_{k \geq n} s_{k,n}. \]

Observe that for $k \geq n+1$,  $s_{k,n} = s_{k,n+1} + (t_{n} - t_{n+1})/2L$. This yields $s_n \leq s_{k,n+1} + (t_n - t_{n+1})/2L$, and thus $s_n \leq s_{n+1} + (t_n - t_{n+1})/2L$, showing that $X_{n+1} \subset X_n$ by \Cref{eq:inclusion}.
Moreover since $s_{k,n}$ is decreasing in $n$, we have $s_{n+1} = \inf_{k \geq n+1} s_{k,n+1} \leq \inf_{k \geq n+1} s_{k,n}$. Since we also have $s_{n+1, n+1} = b_{n+1} \leq b_n = s_{n,n}$, we get $s_{n+1}\leq s_{n}$.

 Since $s_n \leq 1/n$, $s_n$ must converge to 0, so that for $n$ sufficiently large $s_n \leq (c - t_n)/4L$. By construction $s_n \leq S(t_n, \eps)$. From the first point of \Cref{lem:flot_lisse}, on $h_{s_n}^{-1}(t_n + 2s_nL)$ we have $\norme{\nabla h_{s_n}} \geq \rho(1 - \eps)$, so that $X_n$ is a smooth domain. Moreover the bound of the mass of the normal cycle provided by \Cref{prop:sub_dc_aura} applies. Since the Lipschitz constants of $f_{s_n}$ and $g_{s_n}$ are bounded by that of $f$ and $g$ we have
 \[ \limsup_{n \to \infty} M(N_{X_n}) < + \infty.\]
Furthermore by the second point of \Cref{lem:flot_lisse}, since both $t_n$ and $s_n$ converge to zero, $X_n$ converges both acyclically and with respect to the homotopy distance to $X$. 
Finally if a point $x$ belongs to each $X_n$, then $h_{s_n}(x) \leq t_n + 2s_nL$ and letting $n$ tend to $+\infty$ we obtain $h(x) \leq 0$, so that $\cap_n X_n \subset X$. Conversely $X \subset X_n$ for each $n$ as $\norme{h - h_{s_n}}_{\infty} \leq 2s_nL$.
\end{proof}

\end{document}
